% =============================================================================
%  3-Conclusiones.tex — Resultados esperados y requerimientos técnicos
% =============================================================================

\section{Resultados esperados}

Al finalizar el taller, cada participante contará con:

\begin{enumerate}[nosep, leftmargin=1.5em]
  \item Un proyecto QGIS funcional con la red de transporte completa de la ZMVM
    (11~sistemas, más de 11,000~estaciones) conectada en tiempo real a la API de
    Apimetro.
  \item Capas de análisis de cobertura espacial (isócronas de 800~m) y vacíos
    de servicio en las periferias metropolitanas, exportadas en formato
    GeoPackage.
  \item Comprensión práctica del flujo de trabajo
    \textbf{API REST $\to$ QGIS $\to$ análisis $\to$ exportación}, replicable
    en cualquier investigación territorial posterior.
\end{enumerate}

El taller demuestra que la infraestructura de datos abiertos de Apimetro permite
investigación cuantitativa de nivel académico sin requerir que el analista
programe, limpie archivos de texto ni repare trazos geométricos rotos: solo
necesita QGIS y conexión a internet.

\subsection{Requerimientos técnicos del recinto}

\begin{itemize}[nosep, leftmargin=1.5em]
  \item Proyector para el instructor.
  \item Conexiones eléctricas suficientes para las computadoras portátiles de los
    asistentes.
  \item Acceso a la red de internet del recinto, indispensable para realizar las
    peticiones HTTP(S) a la API en tiempo real.
\end{itemize}
